% Storyline 3, paragraph 5: findings and implication.
We find that the direction of the revision rule, rather than the amount of evidence it requires, decides whether private observers agree with one another and still distinguish cooperators from defectors. The beneficial direction depends on the norm, mirroring the kind of misjudgement that private assessment produces under each norm. Compact rules, selected solely by these reputational criteria, make cooperation stable and easy to reach under four classical norms, whereas the same amount of inertia placed on the wrong path can leave no cooperation at all. The rule that best protects cooperators from disputed evidence, however, also shelters opportunists who defect only occasionally, a weakness that the aggregation and score-based rules we test share. Allowing warning states to carry partial consequences restores deterrence while cooperation remains high. Sustaining cooperation through private reputations is therefore not only a matter of what individuals know about each other, but also of how they revise what they believe.
